\documentclass[12pt,letterpaper]{article}

\usepackage[utf8]{inputenc}
\usepackage[T1]{fontenc}
\usepackage[spanish,es-noshorthands]{babel}

\usepackage{amsmath}
\usepackage{amssymb}
\usepackage{amsthm}

\usepackage{geometry}
\usepackage{setspace}
\usepackage{enumitem}
\usepackage{booktabs}
\usepackage{array}
\usepackage{longtable}

\usepackage{xcolor}
\usepackage{titlesec}
\usepackage{fancyhdr}
\usepackage{hyperref}

\geometry{
    top=2.5cm,
    bottom=2.5cm,
    left=2.7cm,
    right=2.7cm
}

\setlength{\parindent}{0pt}
\setlength{\parskip}{0.6em}
\onehalfspacing

\hypersetup{
    colorlinks=true,
    linkcolor=black,
    urlcolor=blue,
    citecolor=black
}

\pagestyle{fancy}
\fancyhf{}
\fancyhead[L]{Métodos Cuantitativos de Investigación en Negocios}
\fancyhead[R]{PhD in Business Economics}
\fancyfoot[C]{\thepage}

\titleformat{\section}
{\large\bfseries}
{\thesection.}
{0.5em}
{}

\newlist{opciones}{enumerate}{1}

\setlist[opciones]{
    label=\arabic*.,
    leftmargin=1.2cm,
    itemsep=0.3em,
    topsep=0.3em
}

\begin{document}

\begin{center}
{\large \textbf{UNIVERSIDAD DEL DESARROLLO}}\\
{\large \textbf{FACULTAD DE ECONOMÍA Y NEGOCIOS}}\\
{\large \textbf{PhD in Business Economics}}

\vspace{1em}

{\Large \textbf{MÉTODOS CUANTITATIVOS DE INVESTIGACIÓN EN NEGOCIOS}}

\vspace{0.7em}

{\large \textbf{TAREA 1: ESCALAS DE MEDICIÓN Y PRUEBAS ESTADÍSTICAS}}
\end{center}

\vspace{1em}

\textbf{Fecha de entrega:} 1 de junio de 2026

\textbf{Ponderación:} 20\%

\vspace{1em}

\section{Enunciado}

\subsection*{1. Escalas de medición}

Para cada una de las siguientes variables, proponga dos escalas de medición (una métrica y una no métrica) de una sola pregunta o ítem. Para esto, debe explicitar qué tipo de escala utiliza (nominal, ordinal, de intervalo o de razón) y describirla completamente en términos de su enunciado y de las opciones de respuesta.

\begin{opciones}
\item Intención de recompra
\item Satisfacción laboral
\item Tamaño de la empresa
\item Credibilidad de mensaje publicitario
\item Antigüedad en la organización
\end{opciones}

\vspace{1em}

\subsection*{2. Constructos reflectivos y formativos}

Identifique un constructo del área de negocios que haya sido medido con una escala reflectiva y uno que haya sido medido con una escala formativa. Para cada uno de ellos, proporcione la definición conceptual del constructo, la escala de medición usada para su operacionalización, y la justificación en relación a su naturaleza (formativa o reflectiva).

\vspace{0.5em}

\textbf{Nota:} Cada estudiante debe proporcionar constructos diferentes, por lo que se recomienda informar previamente los constructos seleccionados para evitar repeticiones.

\vspace{1em}

\subsection*{3. Pruebas estadísticas}

Usando la base ``Datos Tarea 1'' (en Canvas) que contiene los resultados de un experimento en el que se expuso a los participantes a información de una empresa (2 tratamientos con información de actividades de CSR interna o externa de la empresa, y un grupo de control con información general) y se midió su nivel de confianza en la empresa antes y después de la información, se solicita realizar las siguientes pruebas:

\begin{enumerate}[label=\alph*)]
\item Prueba t de medias de dos muestras para la variable edad entre los grupos de tratamiento (no considere el grupo de control).

\item Prueba t de medias para muestras relacionadas para las variables ``Medición de Confianza en la empresa antes de la información'' y ``Medición de Confianza en la empresa después de la información'' para cada uno de los grupos (tratamientos y control).

\item Prueba Z de porcentajes para la variable género entre los grupos de tratamiento (no considere el grupo de control).

\item Prueba de chi-cuadrado para la variable ocupación entre los grupos de tratamiento (no considere el grupo de control).

\item Test ANOVA para las variables ``Medición de Confianza en la empresa antes de la información'' y ``Medición de Confianza en la empresa después de la información'' para cada uno de los grupos (tratamientos y control), incluyendo pruebas post hoc en caso de ser necesario.

\item Test de chi-cuadrado para la variable ingresos en los tres grupos.

\item Test ANOVA para una nueva variable (que usted debe calcular) que mida el cambio en la confianza en la empresa producto de la información entre los tres grupos.
\end{enumerate}

Para cada prueba, reporte los valores del estadístico, grados de libertad (en caso de ser pertinente) y su cálculo, estadístico de comparación al 95\% de confianza, conclusión respecto a la hipótesis nula, e interpretación de los resultados.

\newpage

\section{Respuestas}

\subsection*{1. Escalas de medición}

En investigación cuantitativa, la operacionalización corresponde al proceso mediante el cual un concepto abstracto o constructo teórico es transformado en una variable observable y medible empíricamente. La selección de una escala de medición es particularmente relevante, ya que determina el tipo de información que será recolectada y condiciona posteriormente las técnicas estadísticas que pueden utilizarse en el análisis de datos.

Las escalas nominales permiten clasificar observaciones en categorías mutuamente excluyentes. Las escalas ordinales incorporan además una relación de orden entre categorías, aunque sin asumir distancias equivalentes entre ellas. Las escalas de intervalo permiten interpretar diferencias entre valores y suponen distancias relativamente homogéneas entre categorías. Finalmente, las escalas de razón incorporan adicionalmente un cero absoluto, permitiendo comparaciones de proporción entre observaciones.

En investigación aplicada en negocios, marketing y comportamiento organizacional, muchas variables actitudinales suelen operacionalizarse mediante escalas tipo Likert o escalas numéricas de múltiples categorías. Aunque técnicamente varias de estas escalas poseen naturaleza ordinal, frecuentemente son tratadas empíricamente como aproximaciones métricas debido a que facilitan la aplicación de análisis estadísticos paramétricos y permiten capturar con mayor precisión la intensidad de percepciones, actitudes o intenciones conductuales.

\subsubsection*{a) Intención de recompra}

\textbf{Respuesta}

\textbf{Escala métrica: escala de intervalo}

\textbf{Enunciado}

En una escala de 0 a 10, ¿qué tan probable es que vuelva a comprar productos o servicios de esta empresa durante los próximos seis meses?

\textbf{Opciones de respuesta}

0 = Nada probable

10 = Muy probable

El encuestado puede seleccionar cualquier número entero entre 0 y 10.

\textbf{Justificación metodológica}

La intención de recompra corresponde a una variable actitudinal asociada a una conducta futura esperada del consumidor. En investigación de marketing y comportamiento del consumidor, este tipo de variables suele operacionalizarse mediante escalas de intervalo debido a que interesa capturar no solamente dirección, sino también intensidad de la intención conductual.

La utilización de una escala de 0 a 10 permite obtener mayor variabilidad en las respuestas y discriminar con mayor precisión entre distintos niveles de disposición a recomprar. Aunque técnicamente este tipo de escalas no garantiza propiedades métricas perfectas, en investigación cuantitativa aplicada frecuentemente se asume una aproximación a intervalo cuando existe un número suficiente de categorías de respuesta.

Esta aproximación resulta particularmente útil porque posteriormente permite utilizar técnicas estadísticas paramétricas, tales como análisis de correlación, regresiones lineales, análisis factoriales o modelos estructurales.

\textbf{Escala no métrica: escala ordinal}

\textbf{Enunciado}

¿Cuál de las siguientes alternativas describe mejor su intención de volver a comprar en esta empresa?

\textbf{Opciones de respuesta}

\begin{opciones}
\item Definitivamente no volvería a comprar
\item Probablemente no volvería a comprar
\item No estoy seguro/a
\item Probablemente volvería a comprar
\item Definitivamente volvería a comprar
\end{opciones}

\textbf{Justificación metodológica}

Las categorías poseen un orden lógico creciente respecto de la intención de recompra. Sin embargo, no puede asegurarse que la distancia psicológica entre categorías sea equivalente. Por ejemplo, la diferencia percibida entre ``probablemente volvería'' y ``definitivamente volvería'' puede no ser igual a la diferencia entre otras categorías de respuesta.

Por esta razón, la variable conserva únicamente información jerárquica y debe clasificarse metodológicamente como una escala ordinal.

\subsubsection*{b) Satisfacción laboral}

\textbf{Respuesta}

\textbf{Escala métrica: escala de intervalo}

\textbf{Enunciado}

Considerando su experiencia general en la organización, indique su nivel de satisfacción laboral.

\textbf{Opciones de respuesta}

\begin{opciones}
\item Muy insatisfecho/a
\item Insatisfecho/a
\item Algo insatisfecho/a
\item Ni satisfecho/a ni insatisfecho/a
\item Algo satisfecho/a
\item Satisfecho/a
\item Muy satisfecho/a
\end{opciones}

\textbf{Justificación metodológica}

La satisfacción laboral corresponde a una variable psicológica latente relacionada con percepciones subjetivas del individuo respecto de su experiencia organizacional. Debido a que este tipo de constructos no puede observarse directamente, resulta necesario operacionalizarlo mediante escalas actitudinales.

Aunque las escalas tipo Likert son técnicamente ordinales, en investigación cuantitativa frecuentemente se tratan como aproximaciones métricas cuando poseen múltiples categorías y se asume una distancia relativamente homogénea entre respuestas. Esta práctica es ampliamente utilizada en investigación organizacional, comportamiento organizacional y recursos humanos.

Además, una escala de siete categorías permite capturar con mayor precisión diferencias individuales en la percepción de satisfacción, mejorando la sensibilidad de la medición y facilitando posteriormente la utilización de técnicas estadísticas paramétricas.

\textbf{Escala no métrica: escala ordinal}

\textbf{Enunciado}

¿Cómo clasificaría su nivel actual de satisfacción laboral?

\textbf{Opciones de respuesta}

\begin{opciones}
\item Baja
\item Media
\item Alta
\end{opciones}

\textbf{Justificación metodológica}

La variable presenta una estructura jerárquica natural, ya que ``alta'' representa un mayor nivel de satisfacción que ``media'' o ``baja''. Sin embargo, no es posible conocer la magnitud exacta de la diferencia entre categorías.

Por ello, la escala corresponde metodológicamente a una variable ordinal.

\subsubsection*{c) Tamaño de la empresa}

\textbf{Respuesta}

\textbf{Escala métrica: escala de razón}

\textbf{Enunciado}

¿Cuántos trabajadores tiene actualmente su empresa?

\textbf{Opciones de respuesta}

Respuesta numérica abierta.

Por ejemplo: 5 trabajadores, 40 trabajadores, 250 trabajadores o 1200 trabajadores.

\textbf{Justificación metodológica}

La variable corresponde a una escala de razón debido a que posee cero absoluto y permite realizar comparaciones de magnitud y proporción entre observaciones.

Por ejemplo, una empresa con 200 trabajadores efectivamente posee el doble de trabajadores que una empresa con 100. Esto implica que las diferencias y proporciones entre valores tienen interpretación empírica directa.

Además, este tipo de operacionalización permite utilizar plenamente herramientas estadísticas cuantitativas, tales como análisis descriptivos, modelos econométricos y procedimientos inferenciales paramétricos.

\textbf{Escala no métrica: escala ordinal}

\textbf{Enunciado}

Según el número de trabajadores, ¿en qué categoría se ubica su empresa?

\textbf{Opciones de respuesta}

\begin{opciones}
\item Microempresa: 1 a 9 trabajadores
\item Pequeña empresa: 10 a 49 trabajadores
\item Mediana empresa: 50 a 199 trabajadores
\item Gran empresa: 200 o más trabajadores
\end{opciones}

\textbf{Justificación metodológica}

Las categorías representan distintos niveles crecientes de tamaño organizacional y poseen un orden natural. Sin embargo, las amplitudes entre categorías no son equivalentes, ya que los intervalos contienen rangos de tamaño diferentes.

Por ello, la variable pierde propiedades métricas completas y debe clasificarse metodológicamente como ordinal.

\subsubsection*{d) Credibilidad del mensaje publicitario}

\textbf{Respuesta}

\textbf{Escala métrica: escala de intervalo}

\textbf{Enunciado}

Después de observar el mensaje publicitario, indique qué tan creíble le parece.

\textbf{Opciones de respuesta}

\begin{opciones}
\item Nada creíble
\item Muy poco creíble
\item Poco creíble
\item Medianamente creíble
\item Algo creíble
\item Bastante creíble
\item Muy creíble
\end{opciones}

\textbf{Justificación metodológica}

La credibilidad publicitaria corresponde a una percepción subjetiva del consumidor respecto del grado de confianza y veracidad atribuido al mensaje observado. En investigación de marketing y publicidad, este tipo de percepciones suele operacionalizarse mediante escalas de múltiples categorías para capturar distintos niveles de intensidad actitudinal.

Aunque formalmente la escala posee naturaleza ordinal, en investigación aplicada frecuentemente se aproxima a propiedades de intervalo debido a la gran cantidad de categorías disponibles y a la necesidad de aplicar técnicas estadísticas paramétricas.

Además, utilizar múltiples niveles de respuesta permite aumentar la sensibilidad de la medición y captar diferencias más finas entre consumidores.

\textbf{Escala no métrica: escala nominal}

\textbf{Enunciado}

¿Cómo clasificaría el mensaje publicitario observado?

\textbf{Opciones de respuesta}

\begin{opciones}
\item Creíble
\item No creíble
\item No puedo determinarlo
\end{opciones}

\textbf{Justificación metodológica}

Las categorías clasifican distintas percepciones respecto del mensaje publicitario, pero no establecen necesariamente un orden jerárquico completo entre todas las alternativas.

Por esta razón, la variable corresponde a una escala nominal.

\subsubsection*{e) Antigüedad en la organización}

\textbf{Respuesta}

\textbf{Escala métrica: escala de razón}

\textbf{Enunciado}

¿Cuántos años completos lleva trabajando en esta organización?

\textbf{Opciones de respuesta}

Respuesta numérica abierta. Por ejemplo: 0 años, 3 años, 8 años o 15 años.

\textbf{Justificación metodológica}

La antigüedad laboral corresponde a una variable cuantitativa con cero absoluto, ya que un valor igual a cero representa efectivamente ausencia de antigüedad dentro de la organización.

Esto permite realizar comparaciones de magnitud y proporción entre observaciones. Por ejemplo, un trabajador con 10 años de antigüedad posee el doble de permanencia organizacional que un trabajador con 5 años. En consecuencia, tanto las diferencias como las proporciones entre valores tienen interpretación empírica directa.

\textbf{Escala no métrica: escala ordinal}

\textbf{Enunciado}

¿En qué tramo se encuentra su antigüedad en la organización?

\textbf{Opciones de respuesta}

\begin{opciones}
\item Menos de 1 año
\item Entre 1 y 3 años
\item Entre 4 y 7 años
\item Entre 8 y 15 años
\item Más de 15 años
\end{opciones}

\textbf{Justificación metodológica}

La variable presenta una estructura jerárquica natural, ya que los distintos tramos representan niveles crecientes de permanencia dentro de la organización. Sin embargo, las distancias entre categorías no son equivalentes. Debido a ello, la escala conserva información ordinal, pero pierde propiedades métricas completas.

\subsection*{2. Constructos reflectivos y formativos}

En investigación en negocios, numerosos fenómenos relevantes no pueden observarse directamente. Conceptos como satisfacción, compromiso, confianza, reputación o responsabilidad social corresponden a variables latentes o constructos teóricos que requieren ser operacionalizados mediante indicadores observables.

Dependiendo de la relación causal entre el constructo y sus indicadores, la literatura distingue principalmente entre modelos reflectivos y modelos formativos.

\subsubsection*{Constructo reflectivo: satisfacción del cliente}

\textbf{Respuesta}

\textbf{Definición conceptual}

La satisfacción del cliente corresponde a la evaluación global que realiza un consumidor respecto de su experiencia con una empresa, producto o servicio, considerando el grado en que el desempeño percibido cumple o supera sus expectativas previas. Un ejemplo reciente de medición de este constructo se encuentra en Manhas, Sharma y Quintela (2024), quienes estudian la relación entre innovación de producto, experiencia del cliente y satisfacción del cliente en restaurantes de comida rápida mediante modelos de ecuaciones estructurales.

\textbf{Escala de medición}

En el estudio de Manhas, Sharma y Quintela (2024), la satisfacción del cliente se mide mediante una escala tipo Likert aplicada a clientes de restaurantes de comida rápida. Los indicadores reportados corresponden a afirmaciones sobre satisfacción, agrado, eficiencia del servicio, experiencia de consumo, elección del restaurante, cumplimiento de expectativas y satisfacción global.

\textbf{Indicadores utilizados en el paper}

\begin{opciones}
\item CS1: ``I am extremely pleased with the food and services offered at the QSR''.
\item CS3: ``I am delighted to visit the QSR''.
\item CS4: ``QSR staff offer services effectively and efficiently''.
\item CS6: ``I was happy with the dining experience in the QSR''.
\item CS7: ``I am satisfied with my decision to choose this QSR''.
\item CS8: ``Overall, the QSR meets all my expectations''.
\item CS9: ``QSR gives me overall satisfaction''.
\end{opciones}

Estos indicadores fueron evaluados empíricamente mediante cargas estandarizadas, confiabilidad compuesta y varianza media extraída, lo que es consistente con una lógica de medición reflectiva.

\textbf{Justificación de su naturaleza reflectiva}

En este caso, la satisfacción global del cliente constituye el fenómeno latente subyacente que genera las respuestas observadas en los distintos indicadores. Es decir, los ítems representan manifestaciones observables de una misma evaluación general del consumidor respecto de su experiencia con el restaurante.

Por esta razón, se espera que los indicadores presenten correlaciones relativamente elevadas entre sí. Un individuo altamente satisfecho debería tender a responder positivamente en la mayoría de los ítems de la escala.

Además, los indicadores son relativamente intercambiables: eliminar un ítem puede reducir la precisión o confiabilidad de la medición, pero no cambia sustancialmente el significado conceptual del constructo. Por ello, satisfacción del cliente se interpreta como un constructo reflectivo.

\subsubsection*{Constructo formativo: imagen percibida de responsabilidad social corporativa}

\textbf{Respuesta}

\textbf{Definición conceptual}

La imagen percibida de responsabilidad social corporativa corresponde a la evaluación que realizan los stakeholders respecto del comportamiento social, ambiental, ético y comunitario de una organización. Wang, Hu y Zhang (2020) desarrollan y prueban empíricamente una escala de responsabilidad social corporativa en cadenas hoteleras internacionales, conceptualizando la CSR como un constructo formado por distintas dimensiones, tales como protección ambiental, bienestar de empleados, ética empresarial y bienestar de clientes.

\textbf{Escala de medición}

Escala tipo Likert de 1 a 7, donde:

1 = Totalmente en desacuerdo

7 = Totalmente de acuerdo

\textbf{Dimensiones utilizadas en el paper}

\begin{opciones}
\item Environment protection: prácticas orientadas a la protección ambiental.
\item Employee wellness: prácticas orientadas al bienestar de los empleados.
\item Business ethics: prácticas asociadas a ética empresarial.
\item Customer wellness: prácticas orientadas al bienestar de los clientes.
\end{opciones}

\textbf{Justificación de su naturaleza formativa}

En este caso, los indicadores no representan manifestaciones de un único fenómeno latente subyacente, sino que construyen conjuntamente el constructo global de responsabilidad social corporativa.

Cada dimensión representa una faceta distinta del fenómeno analizado. Por ejemplo, una empresa podría exhibir un desempeño ambiental elevado y simultáneamente presentar debilidades en prácticas laborales, ética empresarial o bienestar de clientes.

Por ello, en modelos formativos los indicadores no necesariamente deben estar altamente correlacionados entre sí. Además, la eliminación de una dimensión modifica sustancialmente el significado conceptual del constructo, ya que implica perder parte de la información que define la responsabilidad social corporativa.

\textbf{Referencias utilizadas}

Manhas, P. S., Sharma, P., \& Quintela, J. A. (2024). Product Innovation and Customer Experience: Catalysts for Enhancing Satisfaction in Quick Service Restaurants. \textit{Tourism and Hospitality, 5}(3), 559--576.

Wang, C., Hu, R., \& Zhang, T. (2020). Corporate social responsibility in international hotel chains and its effects on local employees: Scale development and empirical testing in China. \textit{International Journal of Hospitality Management, 90}, 102598.


\subsection*{3. Pruebas estadísticas}

Para responder esta sección se utilizó la base \textit{Datos Tarea 1}, que contiene 120 observaciones distribuidas equitativamente en tres grupos experimentales: CSR Externa, CSR Interna e Información sin CSR. Cada grupo contiene 40 participantes. Las variables categóricas fueron interpretadas a partir de las etiquetas contenidas en la metadata del archivo SPSS.

Antes de realizar las pruebas inferenciales, se calcularon estadísticos descriptivos generales. La edad promedio fue 30.675 años en CSR Externa, 31.225 años en CSR Interna y 31.900 años en el grupo Información sin CSR. La confianza promedio antes de la información fue 4.700 en CSR Externa, 5.125 en CSR Interna y 4.650 en Información sin CSR. Después de la información, la confianza promedio fue 7.275 en CSR Externa, 6.550 en CSR Interna y 4.900 en Información sin CSR.

\subsubsection*{a) Prueba t de medias de dos muestras para edad entre los grupos de tratamiento}

El objetivo de esta prueba es evaluar si existen diferencias significativas en la edad promedio entre los dos grupos de tratamiento, CSR Externa y CSR Interna. No se considera el grupo de control.

\textbf{Hipótesis nula:}
\[
H_0: \mu_{\text{Edad CSR Externa}} = \mu_{\text{Edad CSR Interna}}
\]

\textbf{Hipótesis alternativa:}
\[
H_1: \mu_{\text{Edad CSR Externa}} \neq \mu_{\text{Edad CSR Interna}}
\]

Se utilizó una prueba t de dos muestras independientes con corrección de Welch, dado que esta versión no exige asumir igualdad de varianzas entre los grupos.

Los resultados descriptivos fueron los siguientes:

\begin{itemize}
    \item CSR Externa: \(n = 40\), media = 30.6750, desviación estándar = 8.3309.
    \item CSR Interna: \(n = 40\), media = 31.2250, desviación estándar = 8.9743.
\end{itemize}

La corrección de Welch se utiliza porque ajusta el error estándar y los grados de libertad cuando las varianzas muestrales de los dos grupos no se asumen iguales. En este caso, en lugar de usar una varianza combinada, se calcula el estadístico t con las varianzas separadas de cada grupo:

\[
t =
\frac{\bar{X}_1 - \bar{X}_2}
{\sqrt{\frac{s_1^2}{n_1} + \frac{s_2^2}{n_2}}}
\]

donde \(\bar{X}_1\) y \(\bar{X}_2\) son las medias de edad de los grupos, \(s_1^2\) y \(s_2^2\) son las varianzas muestrales, y \(n_1\) y \(n_2\) son los tamaños muestrales.

Reemplazando los valores:

\[
t =
\frac{30.6750 - 31.2250}
{\sqrt{\frac{8.3309^2}{40} + \frac{8.9743^2}{40}}}
\]

\[
t =
\frac{-0.5500}
{\sqrt{1.7351 + 2.0134}}
\]

\[
t =
\frac{-0.5500}{1.9361}
= -0.2841
\]

Por lo tanto, el estadístico obtenido fue:

\[
t = -0.2841
\]

Los grados de libertad de Welch se calculan con la siguiente fórmula:

\[
gl =
\frac{
\left(\frac{s_1^2}{n_1} + \frac{s_2^2}{n_2}\right)^2
}{
\frac{\left(\frac{s_1^2}{n_1}\right)^2}{n_1 - 1}
+
\frac{\left(\frac{s_2^2}{n_2}\right)^2}{n_2 - 1}
}
\]

Reemplazando los valores:

\[
gl =
\frac{
\left(\frac{8.3309^2}{40} + \frac{8.9743^2}{40}\right)^2
}{
\frac{\left(\frac{8.3309^2}{40}\right)^2}{40 - 1}
+
\frac{\left(\frac{8.9743^2}{40}\right)^2}{40 - 1}
}
\]

\[
gl = 77.5723
\]

Por lo tanto, los grados de libertad aproximados de Welch fueron:

\[
gl = 77.5723
\]

El estadístico crítico se obtiene desde la distribución t de Student usando una prueba bilateral con \(\alpha = 0.05\). Como la prueba es bilateral, se reparte el nivel de significancia en dos colas:

\[
\alpha/2 = 0.025
\]

Por lo tanto, se busca el percentil:

\[
t_{1-\alpha/2, gl} = t_{0.975, 77.5723}
\]

Esto entrega:

\[
t_c = \pm 1.9910
\]

El p-value fue:

\[
p = 0.7771
\]

Dado que \(p > 0.05\), no se rechaza la hipótesis nula. Por lo tanto, no existe evidencia estadísticamente significativa para afirmar que la edad promedio difiera entre los grupos CSR Externa y CSR Interna. Esto sugiere que ambos grupos de tratamiento son comparables en términos de edad.

\subsubsection*{b) Prueba t de medias para muestras relacionadas para confianza antes y después}

El objetivo de esta prueba es evaluar si la confianza en la empresa cambió significativamente después de recibir la información dentro de cada grupo experimental. Como se compara la medición antes y después para los mismos participantes, corresponde aplicar una prueba t para muestras relacionadas.

Para cada grupo, las hipótesis son:

\textbf{Hipótesis nula:}
\[
H_0: \mu_{\text{Confianza después} - \text{Confianza antes}} = 0
\]

\textbf{Hipótesis alternativa:}
\[
H_1: \mu_{\text{Confianza después} - \text{Confianza antes}} \neq 0
\]

Los grados de libertad para cada prueba se calculan como:

\[
gl = n - 1 = 40 - 1 = 39
\]

El estadístico t pareado se calcula con la media de las diferencias individuales entre confianza después y confianza antes:

\[
t =
\frac{\bar{d}}
{s_d/\sqrt{n}}
\]

donde \(\bar{d}\) es el cambio promedio, \(s_d\) es la desviación estándar de los cambios individuales y \(n\) es el número de pares observados.

El estadístico crítico bilateral se obtiene desde la distribución t de Student con \(gl = 39\). Como la prueba es bilateral y \(\alpha = 0.05\), se busca:

\[
t_{1-\alpha/2,gl} = t_{0.975,39}
\]

Por lo tanto, el estadístico crítico bilateral al 95\% de confianza fue:

\[
t_c = \pm 2.0227
\]

\textbf{Grupo CSR Externa.}

En este grupo, la confianza promedio antes de la información fue 4.7000 y la confianza promedio después fue 7.2750. El cambio promedio fue 2.5750.

La desviación estándar del cambio fue 1.0099. Por lo tanto:

\[
t =
\frac{2.5750}
{1.0099/\sqrt{40}}
=
\frac{2.5750}{0.1597}
= 16.1263
\]

Los resultados de la prueba fueron:

\[
t = 16.1263,\quad gl = 39,\quad p < 0.001
\]

Dado que \(p < 0.05\), se rechaza la hipótesis nula. Por lo tanto, existe evidencia estadísticamente significativa de que la confianza aumentó después de recibir información sobre CSR Externa.

\textbf{Grupo CSR Interna.}

En este grupo, la confianza promedio antes de la información fue 5.1250 y la confianza promedio después fue 6.5500. El cambio promedio fue 1.4250.

La desviación estándar del cambio fue 0.5006. Por lo tanto:

\[
t =
\frac{1.4250}
{0.5006/\sqrt{40}}
=
\frac{1.4250}{0.0792}
= 18.0019
\]

Los resultados de la prueba fueron:

\[
t = 18.0019,\quad gl = 39,\quad p < 0.001
\]

Dado que \(p < 0.05\), se rechaza la hipótesis nula. Por lo tanto, existe evidencia estadísticamente significativa de que la confianza aumentó después de recibir información sobre CSR Interna.

\textbf{Grupo Información sin CSR.}

En este grupo, la confianza promedio antes de la información fue 4.6500 y la confianza promedio después fue 4.9000. El cambio promedio fue 0.2500.

La desviación estándar del cambio fue 1.1491. Por lo tanto:

\[
t =
\frac{0.2500}
{1.1491/\sqrt{40}}
=
\frac{0.2500}{0.1817}
= 1.3759
\]

Los resultados de la prueba fueron:

\[
t = 1.3759,\quad gl = 39,\quad p = 0.1767
\]

Dado que \(p > 0.05\), no se rechaza la hipótesis nula. Por lo tanto, no existe evidencia estadísticamente significativa de que la confianza haya cambiado en el grupo de control.

En síntesis, la confianza aumentó significativamente en los dos grupos que recibieron información de CSR, pero no aumentó significativamente en el grupo que recibió información general sin CSR.

\subsubsection*{c) Prueba Z de porcentajes para género entre los grupos de tratamiento}

El objetivo de esta prueba es comparar la proporción de participantes de género femenino entre los grupos CSR Externa y CSR Interna. No se considera el grupo de control.

\textbf{Hipótesis nula:}
\[
H_0: p_{\text{Femenino CSR Externa}} = p_{\text{Femenino CSR Interna}}
\]

\textbf{Hipótesis alternativa:}
\[
H_1: p_{\text{Femenino CSR Externa}} \neq p_{\text{Femenino CSR Interna}}
\]

En CSR Externa hubo 18 participantes de género femenino de un total de 40, lo que corresponde a una proporción de 0.4500. En CSR Interna hubo 16 participantes de género femenino de un total de 40, lo que corresponde a una proporción de 0.4000.

La prueba Z para dos proporciones utiliza la proporción agrupada bajo la hipótesis nula:

\[
\hat{p} =
\frac{x_1 + x_2}{n_1 + n_2}
\]

Reemplazando los valores:

\[
\hat{p} =
\frac{18 + 16}{40 + 40}
=
\frac{34}{80}
= 0.4250
\]

El error estándar se calcula como:

\[
EE =
\sqrt{\hat{p}(1-\hat{p})\left(\frac{1}{n_1}+\frac{1}{n_2}\right)}
\]

\[
EE =
\sqrt{0.4250(1-0.4250)\left(\frac{1}{40}+\frac{1}{40}\right)}
= 0.1105
\]

Luego, el estadístico Z se calcula como:

\[
Z =
\frac{\hat{p}_1 - \hat{p}_2}{EE}
=
\frac{0.4500 - 0.4000}{0.1105}
= 0.4523
\]

Por lo tanto, el estadístico obtenido fue:

\[
Z = 0.4523
\]

El estadístico crítico se obtiene desde la distribución normal estándar. Como la prueba es bilateral y \(\alpha = 0.05\), se busca:

\[
Z_{1-\alpha/2} = Z_{0.975}
\]

Por lo tanto, el estadístico crítico bilateral al 95\% de confianza fue:

\[
Z_c = \pm 1.9600
\]

El p-value fue:

\[
p = 0.6510
\]

Dado que \(p > 0.05\), no se rechaza la hipótesis nula. Por lo tanto, no existe evidencia estadísticamente significativa para afirmar que la proporción de mujeres difiera entre los grupos CSR Externa y CSR Interna. Esto sugiere que ambos grupos de tratamiento son comparables en términos de género.

\subsubsection*{d) Prueba de chi-cuadrado para ocupación entre los grupos de tratamiento}

El objetivo de esta prueba es evaluar si la distribución de la variable ocupación difiere entre los grupos CSR Externa y CSR Interna. No se considera el grupo de control.

\textbf{Hipótesis nula:}
\[
H_0: \text{la ocupación es independiente del grupo de tratamiento}
\]

\textbf{Hipótesis alternativa:}
\[
H_1: \text{la ocupación no es independiente del grupo de tratamiento}
\]

Se utilizó una prueba chi-cuadrado de independencia. La tabla observada fue:

\[
\begin{array}{lrrrr}
\hline
 & \text{Estudiante} & \text{Otro} & \text{Trab. dep.} & \text{Trab. indep.} \\
\hline
\text{CSR Externa} & 7 & 14 & 7 & 12 \\
\text{CSR Interna} & 9 & 9 & 12 & 10 \\
\hline
\end{array}
\]

Las frecuencias esperadas se calculan como:

\[
E_{ij} =
\frac{(\text{total fila}_i)(\text{total columna}_j)}
{\text{total general}}
\]

Por ejemplo, para estudiantes en CSR Externa:

\[
E =
\frac{40 \times 16}{80}
= 8.0000
\]

La tabla esperada fue:

\[
\begin{array}{lrrrr}
\hline
 & \text{Estudiante} & \text{Otro} & \text{Trab. dep.} & \text{Trab. indep.} \\
\hline
\text{CSR Externa} & 8.0 & 11.5 & 9.5 & 11.0 \\
\text{CSR Interna} & 8.0 & 11.5 & 9.5 & 11.0 \\
\hline
\end{array}
\]

El estadístico chi-cuadrado se calcula como:

\[
\chi^2 =
\sum \frac{(O_{ij}-E_{ij})^2}{E_{ij}}
\]

Al reemplazar las frecuencias observadas y esperadas de la tabla, se obtiene:

\[
\chi^2 = 2.8346
\]

Los grados de libertad se calcularon como:

\[
gl = (r-1)(c-1) = (2-1)(4-1) = 3
\]

El estadístico crítico se obtiene desde la distribución chi-cuadrado con \(gl = 3\):

\[
\chi^2_{1-\alpha,gl} = \chi^2_{0.95,3}
\]

Por lo tanto, el estadístico crítico al 95\% de confianza fue:

\[
\chi^2_c = 7.8147
\]

El p-value fue:

\[
p = 0.4178
\]

Dado que \(p > 0.05\), no se rechaza la hipótesis nula. Por lo tanto, no existe evidencia estadísticamente significativa de asociación entre ocupación y grupo de tratamiento. Esto sugiere que los grupos CSR Externa y CSR Interna son comparables en términos de ocupación.

\subsubsection*{e) ANOVA para confianza antes y después entre los tres grupos}

El objetivo de esta prueba es comparar las medias de confianza entre los tres grupos experimentales: CSR Externa, CSR Interna e Información sin CSR. Se realizaron dos ANOVA de una vía: uno para la confianza antes de la información y otro para la confianza después de la información.

Para cada ANOVA, las hipótesis son:

\textbf{Hipótesis nula:}
\[
H_0: \mu_1 = \mu_2 = \mu_3
\]

\textbf{Hipótesis alternativa:}
\[
H_1: \text{al menos una media de grupo es diferente}
\]

Los grados de libertad se calcularon como:

\[
gl_{\text{entre}} = k - 1 = 3 - 1 = 2
\]

\[
gl_{\text{dentro}} = N - k = 120 - 3 = 117
\]

El estadístico ANOVA se calcula como la razón entre el cuadrado medio entre grupos y el cuadrado medio dentro de los grupos:

\[
F =
\frac{CM_{\text{entre}}}{CM_{\text{dentro}}}
\]

donde:

\[
CM_{\text{entre}} =
\frac{SC_{\text{entre}}}{gl_{\text{entre}}}
\]

\[
CM_{\text{dentro}} =
\frac{SC_{\text{dentro}}}{gl_{\text{dentro}}}
\]

El estadístico crítico se obtiene desde la distribución F con \(gl_{\text{entre}} = 2\) y \(gl_{\text{dentro}} = 117\):

\[
F_{1-\alpha,gl_{\text{entre}},gl_{\text{dentro}}}
=
F_{0.95,2,117}
= 3.0738
\]

\textbf{ANOVA para confianza antes.}

Las medias de confianza antes de la información fueron:

\begin{itemize}
    \item CSR Externa: media = 4.7000, desviación estándar = 1.2850.
    \item CSR Interna: media = 5.1250, desviación estándar = 1.3623.
    \item Información sin CSR: media = 4.6500, desviación estándar = 1.0990.
\end{itemize}

El resultado del ANOVA fue:

\[
SC_{\text{entre}} = 5.4500,\quad
SC_{\text{dentro}} = 183.8750
\]

\[
CM_{\text{entre}} =
\frac{5.4500}{2}
= 2.7250
\]

\[
CM_{\text{dentro}} =
\frac{183.8750}{117}
= 1.5716
\]

\[
F =
\frac{2.7250}{1.5716}
= 1.7339
\]

\[
F = 1.7339,\quad gl = (2,117),\quad F_c = 3.0738,\quad p = 0.1811
\]

Dado que \(p > 0.05\), no se rechaza la hipótesis nula. Por lo tanto, antes de la intervención no existen diferencias estadísticamente significativas en confianza entre los tres grupos. Esto sugiere que los grupos partían desde niveles similares de confianza.

\textbf{ANOVA para confianza después.}

Las medias de confianza después de la información fueron:

\begin{itemize}
    \item CSR Externa: media = 7.2750, desviación estándar = 1.5357.
    \item CSR Interna: media = 6.5500, desviación estándar = 1.3578.
    \item Información sin CSR: media = 4.9000, desviación estándar = 1.3359.
\end{itemize}

El resultado del ANOVA fue:

\[
SC_{\text{entre}} = 118.5167,\quad
SC_{\text{dentro}} = 233.4750
\]

\[
CM_{\text{entre}} =
\frac{118.5167}{2}
= 59.2583
\]

\[
CM_{\text{dentro}} =
\frac{233.4750}{117}
= 1.9955
\]

\[
F =
\frac{59.2583}{1.9955}
= 29.6958
\]

\[
F = 29.6958,\quad gl = (2,117),\quad F_c = 3.0738,\quad p < 0.001
\]

Dado que \(p < 0.05\), se rechaza la hipótesis nula. Por lo tanto, después de la información existen diferencias estadísticamente significativas en confianza entre los grupos.

Como el ANOVA fue significativo, se aplicó una prueba post hoc Tukey HSD. Los resultados indican que:

\begin{itemize}
    \item CSR Externa difiere significativamente del grupo Información sin CSR.
    \item CSR Interna difiere significativamente del grupo Información sin CSR.
    \item CSR Externa y CSR Interna no difieren significativamente entre sí al 5\%.
\end{itemize}

En términos sustantivos, ambos grupos que recibieron información sobre CSR presentan mayor confianza posterior que el grupo de control. Sin embargo, la diferencia entre CSR Externa y CSR Interna no alcanza significancia estadística al 5\%.

\subsubsection*{f) Test de chi-cuadrado para ingresos en los tres grupos}

El objetivo de esta prueba es evaluar si la distribución de ingresos difiere entre los tres grupos experimentales.

\textbf{Hipótesis nula:}
\[
H_0: \text{los ingresos son independientes del grupo experimental}
\]

\textbf{Hipótesis alternativa:}
\[
H_1: \text{los ingresos no son independientes del grupo experimental}
\]

Se aplicó una prueba chi-cuadrado de independencia. La tabla observada fue:

\[
\begin{array}{lrrrrr}
\hline
 & \text{1 a 2 MM} & \text{200 a 500 mil} & \text{501 a 1000 mil} & \text{<200 mil} & \text{>2 MM} \\
\hline
\text{CSR Externa} & 9 & 10 & 6 & 10 & 5 \\
\text{CSR Interna} & 11 & 7 & 15 & 3 & 4 \\
\text{Info sin CSR} & 11 & 9 & 5 & 10 & 5 \\
\hline
\end{array}
\]

Las frecuencias esperadas se calculan como:

\[
E_{ij} =
\frac{(\text{total fila}_i)(\text{total columna}_j)}
{\text{total general}}
\]

Por ejemplo, para el tramo de ingresos entre 1 millón y 2 millones en CSR Externa:

\[
E =
\frac{40 \times 31}{120}
= 10.3333
\]

La tabla esperada fue:

\[
\begin{array}{lrrrrr}
\hline
 & \text{1 a 2 MM} & \text{200 a 500 mil} & \text{501 a 1000 mil} & \text{<200 mil} & \text{>2 MM} \\
\hline
\text{CSR Externa} & 10.3333 & 8.6667 & 8.6667 & 7.6667 & 4.6667 \\
\text{CSR Interna} & 10.3333 & 8.6667 & 8.6667 & 7.6667 & 4.6667 \\
\text{Info sin CSR} & 10.3333 & 8.6667 & 8.6667 & 7.6667 & 4.6667 \\
\hline
\end{array}
\]

El estadístico chi-cuadrado se calcula como:

\[
\chi^2 =
\sum \frac{(O_{ij}-E_{ij})^2}{E_{ij}}
\]

Al reemplazar las frecuencias observadas y esperadas de la tabla, se obtiene:

\[
\chi^2 = 12.2003
\]

Los grados de libertad se calcularon como:

\[
gl = (r-1)(c-1) = (3-1)(5-1) = 8
\]

El estadístico crítico se obtiene desde la distribución chi-cuadrado con \(gl = 8\):

\[
\chi^2_{1-\alpha,gl} = \chi^2_{0.95,8}
\]

Por lo tanto, el estadístico crítico al 95\% de confianza fue:

\[
\chi^2_c = 15.5073
\]

El p-value fue:

\[
p = 0.1425
\]

Dado que \(p > 0.05\), no se rechaza la hipótesis nula. Por lo tanto, no existe evidencia estadísticamente significativa para afirmar que la distribución de ingresos difiera entre los tres grupos. Esto sugiere que los grupos son comparables en términos de ingresos.

\subsubsection*{g) ANOVA para el cambio en confianza entre los tres grupos}

Para esta prueba se construyó una nueva variable denominada cambio en confianza, calculada como:

\[
\text{Cambio en confianza} = \text{Confianza después} - \text{Confianza antes}
\]

El objetivo es evaluar si el cambio promedio en confianza difiere entre los tres grupos.

\textbf{Hipótesis nula:}
\[
H_0: \mu_{\Delta 1} = \mu_{\Delta 2} = \mu_{\Delta 3}
\]

\textbf{Hipótesis alternativa:}
\[
H_1: \text{al menos un grupo tiene un cambio promedio diferente}
\]

Las medias de cambio fueron:

\begin{itemize}
    \item CSR Externa: media = 2.5750, desviación estándar = 1.0099.
    \item CSR Interna: media = 1.4250, desviación estándar = 0.5006.
    \item Información sin CSR: media = 0.2500, desviación estándar = 1.1491.
\end{itemize}

Los grados de libertad fueron:

\[
gl_{\text{entre}} = 2,\quad gl_{\text{dentro}} = 117
\]

Al igual que en el ANOVA anterior, el estadístico F se calcula como:

\[
F =
\frac{CM_{\text{entre}}}{CM_{\text{dentro}}}
\]

con:

\[
CM_{\text{entre}} =
\frac{SC_{\text{entre}}}{gl_{\text{entre}}}
\]

\[
CM_{\text{dentro}} =
\frac{SC_{\text{dentro}}}{gl_{\text{dentro}}}
\]

Para la variable cambio en confianza, las sumas de cuadrados fueron:

\[
SC_{\text{entre}} = 108.1167,\quad
SC_{\text{dentro}} = 101.0500
\]

Por lo tanto:

\[
CM_{\text{entre}} =
\frac{108.1167}{2}
= 54.0583
\]

\[
CM_{\text{dentro}} =
\frac{101.0500}{117}
= 0.8637
\]

\[
F =
\frac{54.0583}{0.8637}
= 62.5910
\]

El estadístico crítico se obtiene desde la distribución F:

\[
F_{1-\alpha,gl_{\text{entre}},gl_{\text{dentro}}}
=
F_{0.95,2,117}
= 3.0738
\]

El resultado del ANOVA fue:

\[
F = 62.5910,\quad gl = (2,117),\quad F_c = 3.0738,\quad p < 0.001
\]

Dado que \(p < 0.05\), se rechaza la hipótesis nula. Por lo tanto, el cambio promedio en confianza difiere significativamente entre los tres grupos.

Como el ANOVA fue significativo, se aplicó una prueba post hoc Tukey HSD. Los resultados muestran que los tres pares de grupos difieren significativamente entre sí:

\begin{itemize}
    \item CSR Externa difiere significativamente de CSR Interna.
    \item CSR Externa difiere significativamente de Información sin CSR.
    \item CSR Interna difiere significativamente de Información sin CSR.
\end{itemize}

En términos sustantivos, la información sobre CSR Externa produjo el mayor aumento promedio en confianza, seguida por la información sobre CSR Interna. El grupo Información sin CSR presentó el menor cambio promedio. Este resultado sugiere que el tipo de información recibida influye significativamente en el cambio de confianza hacia la empresa.

\subsubsection*{Conclusión general de la pregunta 3}

En conjunto, los resultados muestran que los grupos son comparables en variables sociodemográficas relevantes, ya que no se encontraron diferencias significativas en edad, género, ocupación ni ingresos. Además, antes de recibir la información, los tres grupos no presentaban diferencias significativas en confianza hacia la empresa.

Después de la intervención, sí se observaron diferencias significativas en confianza entre los grupos. Los grupos expuestos a información sobre CSR mostraron mayores niveles de confianza que el grupo de control. Además, el análisis del cambio en confianza indica que la CSR Externa generó el mayor aumento promedio, seguida por la CSR Interna, mientras que la información general sin CSR produjo un cambio pequeño y no significativo.

Por lo tanto, la evidencia sugiere que la información sobre responsabilidad social corporativa aumenta la confianza en la empresa, y que el efecto es especialmente fuerte cuando la información se refiere a actividades de CSR Externa.


\end{document}
